\bookchapter{פרק 7: סקירה מקיפה של פונקציות מתקדמות}

\chapterintro{הטכנולוגיה מתקדמת במהירות, ופונקציות הטלפון מתפתחות ללא הרף. תקשורת לוויינית, תקשורת ללא רשת, זיהוי תאונות דרכים, טכנולוגיית Star Flash... פונקציות שנשמעות כמו מדע בדיוני כבר נכנסו לחיינו. בואו נכיר את הטכנולוגיות המתקדמות האלו.}

\booksection{תקשורת לוויינית: קריאה לעזרה גם בלי קליטה}

הטלפון יוצר חיבור ישיר עם לווייני תקשורת במסלול ושולח הודעות מצוקה. נכון לעכשיו, הפונקציה תומכת בעיקר בשליחת הודעות טקסט ומידע מיקום, לא ניתן לבצע שיחות.

\booksubsection{עקרון הפעולה：}

Apple

\booksubsection{דגמים נתמכים：}

\begin{itemize}
  \item \textbf{אנדרואיד}iPhone 14 ואילך תומך בשליחת SOS חירום באמצעות לוויין. סדרת iPhone 17 תומכת גם בשליחת iMessage ו-SMS דרך לוויין, והרחיבה כיסוי לאזורים נוספים.
  \item \textbf{OnePlus}:החל משנת 2024, דגמי הדגל תומכים בהדרגה בפונקציית הודעות לוויין BeiDou.
  \item \textbf{מגמת התפתחות}:תקשורת לוויינית הופכת מתכונה בלעדית לדגמי דגל יוקרתיים לתכונה נפוצה יותר, ובעתיד גם דגמי ביניים צפויים לתמוך בה. מערכת הלוויינים הסינית BeiDou נהנית מהתמיכה הרחבה ביותר בטלפונים מתוצרת סין.
  \item \textbf{תרחישי שימוש:}טיולי שטח, מצוקה בטיפוס הרים.
\end{itemize}

\booksubsection{תרחישי שימוש：}

\begin{itemize}
  \item אסונות טבע הגורמים לניתוק תקשורת קרקעית.
  \item אופן השימוש:
  \item וודא שאתה בשטח פתוח עם קו ראייה לשמיים.
\end{itemize}

\booksubsection{אופן השימוש：}

\begin{enumerate}
  \item עקוב אחר ההנחיות במסך וכוון לכיוון הלוויין.
  \item שלח הודעת מצוקה מוגדרת מראש ומיקום.
  \item טיפ קטן
  \item תקשורת לוויינית דורשת אנטנת חומרה מיוחדת בטלפון, לא כל הטלפונים תומכים בכך. בנוסף, תקשורת לוויינית בדרך כלל דורשת מנוי בתשלום, אך פונקציית קריאת החירום בדרך כלל חינמית.
\end{enumerate}

\begin{tcbinfo}{2. תקשורת ללא רשת: דרך חדשה לתקשורת מטווח קצר}
“תקשורת ללא רשת” היא טכנולוגיה המאפשרת לשני טלפונים לתקשר ישירות זה עם זה ללא רשת סלולרית או Wi-Fi.

\end{tcbinfo}

\booksection{תקשורת ללא רשת: דרך חדשה לתקשורת מטווח קצר}

שימוש בבלוטות' או בתדרי רדיו ייעודיים ליצירת חיבור נקודה-לנקודה בטווח קצר (בדרך כלל מאות מטרים עד מספר קילומטרים).

\booksubsection{עקרון טכני：}

פעילויות שטח

\booksubsection{תרחישי יישום：}

\begin{itemize}
  \item \textbf{חילוץ באסונות}:כאשר התקשורת מנותקת לאחר רעידת אדמה או שיטפון.
  \item \textbf{אירועים גדולים}:בהופעות ואירועי ספורט כאשר יש עומס אנשים וקליטה גרועה.
  \item \textbf{פלטפורמות נתמכות:}Apple
\end{itemize}

\booksubsection{פלטפורמות נתמכות：}

\begin{itemize}
  \item \textbf{אנדרואיד}:אפליקציית ChangLian תומכת בשיחות ללא רשת.
  \item \textbf{OnePlus}:חלק מהפונקציות של Xiaomi Miaoxiang Center.
  \item \textbf{Huawei}זיהוי תאונות דרכים הוא פונקציה המשתמשת בחיישני הטלפון לזיהוי תאונות דרכים חמורות וחיוג אוטומטי למספרי חירום.
\end{itemize}

\booksection{זיהוי תאונות דרכים: מציל חיים ברגע הקריטי}

הטלפון משתמש במד תאוצה, גירוסקופ, ברומטר, מיקרופון וחיישנים נוספים כדי לקבוע באופן מקיף אם התרחשה התנגשות חמורה. לאחר זיהוי תאונה, הטלפון:

\booksubsection{עקרון הפעולה：}

אם לא הגבת עד סוף הספירה לאחור, מחייג אוטומטית למספרי חירום.

\begin{enumerate}
  \item שולח את המיקום שלך לאנשי קשר לשעת חירום.
  \item דגמים נתמכים:
  \item אפל
\end{enumerate}

\booksubsection{דגמים נתמכים:}

\begin{itemize}
  \item \textbf{אנדרואיד}: חלק ממכשירי הפרימיום נתמכים בהדרגה.
  \item \textbf{HarmonyOS}יש להפעיל את הפונקציה כדי להשתמש בה.
\end{itemize}

\booksubsection{הערות:}

\begin{itemize}
  \item מסתמכת על חיישנים ואלגוריתמים, אינה אמינה ב-100\%, אך מספקת שכבת הגנה נוספת.
  \item 4. טכנולוגיית SparkLink: הדור הבא של קישוריות אלחוטית
  \item SparkLink (SparkLink) היא טכנולוגיית תקשורת אלחוטית מהדור החדש שפותחה בהובלת Huawei, במטרה להחליף Bluetooth ו-Wi-Fi בתרחישים מסוימים. היא משווה ואף עולה על הפתרונות הקיימים במדדים מרכזיים: קצב נתונים באוויר של עד 900Mbps (הרבה מעבר ל-50Mbps של Bluetooth), השהיה נמוכה עד 20 $\mu$ s (Bluetooth כ-15ms), תמיכה בעד 4096 מכשירים מחוברים בו-זמנית ברשת אחת (Bluetooth בדרך כלל עד 8).
\end{itemize}

\booksection{טכנולוגיית SparkLink: הדור הבא של קישוריות אלחוטית}

השהיה נמוכה במיוחד

\booksubsection{מאפיינים טכניים:}

\begin{itemize}
  \item \textbf{קצב גבוה}: קצב שיא של 900Mbps לקישור בודד, גבוה עשרות מונים מ-Bluetooth אודיו, מאפשר העברת אודיו ללא אובדן ואפילו וידאו HD.
  \item \textbf{צריכת חשמל נמוכה}: צריכת המתנה דומה ל-Bluetooth, אך יעילות העברה גבוהה יותר, צריכת אנרגיה נמוכה יותר ליחידת נתונים.
  \item \textbf{כיסוי רחב + ריבוי חיבורים}: מרחק חיבור גדול ב-50\% מ-Bluetooth, תמיכה באלפי מכשירים בו-זמנית ברשת אחת.
  \item \textbf{תרחישי שימוש:}אוזניות אלחוטיות
\end{itemize}

\booksubsection{תרחישי שימוש:}

\begin{itemize}
  \item \textbf{שעון חכם}: חיבור יציב יותר.
  \item \textbf{שעון חכם}: יותר מכשירים מחוברים בו-זמנית.
  \item \textbf{קישוריות לרכב}: חיבור חלק יותר בין טלפון לרכב.
  \item \textbf{מכשירים נתמכים:}נכון לעכשיו, בעיקר מכשירים במערכת האקולוגית של Huawei, כגון סדרת Huawei Mate 60, חלק מאוזניות ושעוני Huawei. בעתיד יצטרפו עוד מותגים.
\end{itemize}

\booksubsection{מכשירים נתמכים:}

טכנולוגיית SparkLink מייצגת את הכיוון העתידי של תקשורת אלחוטית. ככל שיותר יצרנים יצטרפו, נוכל לראות יותר מכשירים תומכי SparkLink וליהנות מחוויית אלחוט טובה יותר.

\textbf{בדוק: האם הטלפון שלך תומך בתכונות אלה}\par

חפש “לוויין” בהגדרות וראה אם יש אפשרויות הקשורות לתקשורת לוויינית.

\begin{tcbpractice}{חפש "זיהוי תאונות דרכים" או "SOS חירום", בדוק אם מופעל.}
\begin{enumerate}
  \item אם הטלפון שלך הוא Huawei, בדוק אם הוא תומך בחיבור התקני SparkLink.
  \item הבן את שיטות השימוש בתכונות אלה, הן עשויות להועיל ברגעים קריטיים.
  \item 2. פרק התוכנה (פרקים 8-14)
  \item לאחר שליטה בפעולות הבסיסיות, בואו נעמיק בטיפים לשימוש באפליקציות שונות. מתקשורת ועד משרד, מתשלומים ועד בידור, פרקים אלה יעזרו לך להפוך למומחה טלפונים חכם אמיתי.
\end{enumerate}

\end{tcbpractice}
